@section Printing
The @code{Printer} class is used for all aspects of printing.
It is a subclass of the Dynace @code{Stream} class and as
such inherits all of its functionality.

All methods in the class use an argument referred to as @code{pntr}.
This will always be the printer object returned by @code{QueryPrinter},
@code{New} or @code{NewWithHDC} and used to identify the printer
which is to be affected.

The value returned by all methods which cause text or graphics output is
very significant.  During normal operation, the printer object passed
will be returned.  However, if an error occurs or the user aborted the
report a @code{NULL} will be returned.  In this case further output
should be avoided (although it won't hurt - it'll just be ignored) and
the printer object should be disposed.

All positioning parameters are in increments dictated by
@code{SetScale}.  Positions begin in the upper left hand corner and have
an index origin of 0.



















@deffn {Arc} Arc::Printer
@sp 2
@example
@group
r = gArc(pntr, yBeg, xBeg, yEnd, xEnd, yaBeg, xaBeg, yaEnd, xaEnd);

object  pntr;   /*  printer object         */
int     yBeg;   /*  starting row           */
int     xBeg;   /*  starting column        */
int     yEnd;   /*  ending row             */
int     xEnd;   /*  ending column          */
int     yaBeg;  /*  starting arc row       */
int     xaBeg;  /*  starting arc column    */
int     yaEnd;  /*  ending arc row         */
int     xaEnd;  /*  ending arc column      */
object  r;      /*  printer object         */
@end group
@end example
This method is used to output an elliptical arc to the printer.  The
currently selected pen object will indicate the thickness and pattern of
the outline of the shape, and the currently selected brush object will
be used to determine what pattern the shape will be filled with.  The
parameters indicate location of the beginning and ending of the shape as
well as the coordinates of the arc and are in increments dictated by
@code{SetScale}.

See the note at the beginning of this section regarding the return value.
@example
@group
@exdent Example:

object  pntr;

if (!gArc(pntr, 10, 10, 180, 140, 12, 12, 128, 135))
        abort report;
@end group
@end example
@sp 1
See also:  @code{NewPage, SetScale, Use,}
        @code{Line, Rectangle, Ellipse, RoundRect,}
@iftex
@hfil @break @hglue .63in 
@end iftex
@code{Chord, Pie}
@end deffn












@deffn {Chord} Chord::Printer
@sp 2
@example
@group
r = gChord(pntr, yBeg, xBeg, yEnd, xEnd,
                 ylBeg, xlBeg, ylEnd, xlEnd);

object  pntr;   /*  printer object         */
int     yBeg;   /*  starting row           */
int     xBeg;   /*  starting column        */
int     yEnd;   /*  ending row             */
int     xEnd;   /*  ending column          */
int     ylBeg;  /*  starting line row      */
int     xlBeg;  /*  starting line column   */
int     ylEnd;  /*  ending line row        */
int     xlEnd;  /*  ending line column     */
object  r;      /*  printer object         */
@end group
@end example
This method is used to output a chord shape to the printer.  The
currently selected pen object will indicate the thickness and pattern of
the outline of the shape, and the currently selected brush object will
be used to determine what pattern the shape will be filled with.  The
parameters indicate location of the beginning and ending of the shape as
well as the coordinates of the intersecting line and are in increments
dictated by @code{SetScale}.

See the note at the beginning of this section regarding the return value.
@example
@group
@exdent Example:

object  pntr;

if (!gChord(pntr, 10, 10, 180, 140, 12, 12, 128, 135))
        abort report;
@end group
@end example
@sp 1
See also:  @code{NewPage, SetScale, Use,}
        @code{Line, Rectangle, Ellipse, RoundRect, Pie, Arc}
@end deffn




@deffn {Copies} Copies::Printer
@sp 2
@example
@group
n = gCopies(prt);

object  prt;  /*  printer object     */
int     n;    /*  number of copies   */
@end group
@end example
This instance method returns the number of copies set for this print job.
@example
@group
@exdent Example:

object  prt;
int     nc;

nc = gCopies(prt);
@end group
@end example
@end deffn








@deffn {DeepDispose} DeepDispose::Printer
@sp 2
This method performs the same function as @code{Dispose}.  See that
method for details.
@end deffn





@deffn {Dispose} Dispose::Printer
@sp 2
@example
@group
r = gDispose(pntr); 

object  pntr;   /*  printer object  */
object  r;      /*  NULL            */
@end group
@end example
This method is used to flush the final output page, close the printer and
dispose of the printer object.  It must be called when a report is complete and the printer is no longer needed.

The value returned is always @code{NULL} and may be used to null out
the variable which contained the object being disposed in order to
avoid future accidental use.
@example
@group
@exdent Example:

object  pntr;

pntr = gDispose(pntr);
@end group
@end example
@c @sp 1
@c See also:  @code{}
@end deffn
















@deffn {Ellipse} Ellipse::Printer
@sp 2
@example
@group
r = gEllipse(pntr, yBeg, xBeg, yEnd, xEnd);

object  pntr;   /*  printer object   */
int     yBeg;   /*  starting row     */
int     xBeg;   /*  starting column  */
int     yEnd;   /*  ending row       */
int     xEnd;   /*  ending column    */
object  r;      /*  printer object   */
@end group
@end example
This method is used to output an ellipse to the printer.  The currently
selected pen object will indicate the thickness and pattern of the
outline of the shape, and the currently selected brush object will be
used to determine what pattern the shape will be filled with.  The
parameters indicate location of the beginning and ending of the shape
and are in increments dictated by @code{SetScale}.

See the note at the beginning of this section regarding the return value.
@example
@group
@exdent Example:

object  pntr;

if (!gEllipse(pntr, 10, 20, 30, 40))
        abort report;
@end group
@end example
@sp 1
See also:  @code{NewPage, SetScale, Use,}
        @code{Line, Rectangle, RoundRect, Chord, Pie, Arc}
@end deffn













@deffn {Handle} Handle::Printer
@sp 2
@example
@group
r = gHandle(pntr); 

object  pntr;   /*  printer object         */
HANDLE  hdc;    /*  handle device context  */
@end group
@end example
This method is used to obtain the device context handle associated with
an opened printer represented by @code{pntr}.  It is used internally
by Windows and should not normally be needed.
@example
@group
@exdent Example:

HDC     hdc;

hdc = (HDC) gHandle(pntr);
@end group
@end example
@c @sp 1
@c See also:  @code{}
@end deffn

















@deffn {Line} Line::Printer
@sp 2
@example
@group
r = gLine(pntr, yBeg, xBeg, yEnd, xEnd);

object  pntr;   /*  printer object   */
int     yBeg;   /*  starting row     */
int     xBeg;   /*  starting column  */
int     yEnd;   /*  ending row       */
int     xEnd;   /*  ending column    */
object  r;      /*  printer object   */
@end group
@end example
This method is used to output a line to the printer.  The currently
selected pen object will indicate the thickness and pattern of the line.
The parameters indicate location of the beginning and ending of the line
and are in increments dictated by @code{SetScale}.

See the note at the beginning of this section regarding the return value.
@example
@group
@exdent Example:

object  pntr;

if (!gLine(pntr, 10, 20, 30, 40))
        abort report;
@end group
@end example
@sp 1
See also:  @code{NewPage, SetScale, Use,}
        @code{Arc, Rectangle, Ellipse, RoundRect, Chord, Pie}
@end deffn












@deffn {LoadFont} LoadFont::Printer
@sp 2
@example
@group
r = gLoadFont(pntr, fname, sz);

object   pntr;  /*  printer object   */
char    *fname; /*  font name        */
int      sz;    /*  point size       */
object   r;     /*  font object      */
@end group
@end example
This method is used to load an arbitrary font by name at any point size
and associate it as the default font for future text output to printer
@code{pntr}.  @code{fname} is the full name of the font as it appears
when you list the available fonts via the control-panel / fonts Windows
utility, minus the font type in parentheses.  @code{sz} indicates the
desired point size.

Any font object previously associated with the printer will be
disposed.  When the printer object is disposed, the font object
will also be disposed.

The value returned is an object representing the font loaded, or
@code{NULL} if the font was not found.

Note that @code{fname} may also be a Dynace object cast as a (@code{char *}).
@example
@group
@exdent Example:

object  pntr;

gLoadFont(pntr, "Times New Roman", 12);
@end group
@end example
@sp 1
See also:  @code{LoadSystemFont, Indirect::ExternalFont, Use}
@end deffn














@deffn {LoadSystemFont} LoadSystemFont::Printer
@sp 2
@example
@group
r = gLoadSystemFont(pntr, fnt);

object   pntr;  /*  a printer object  */
unsigned fnt;   /*  font identifier  */
object   r;     /*  font object      */
@end group
@end example
This method is used to load a Windows predefined font and associate it
with printer @code{pntr}.  The last font associated with a printer is the
one which will be used when any text is output to the printer.

@code{fnt} is a Windows defined macro which identifies the font.  The
available options are defined in the Windows documentation under the
function named @code{GetStockObject} and normally end with @code{_FONT}.
The font object will be automatically destroyed whenever the printer
object is disposed.

The value returned is an object representing the font, or
@code{NULL} if the font was not found.
@example
@group
@exdent Example:

object  pntr;

gLoadSystemFont(pntr, SYSTEM_FONT);
@end group
@end example
@sp 1
See also:  @code{LoadFont, Load::SystemFont, Use}
@end deffn

















@deffn {New} New::Printer
@sp 2
@example
@group
pntr = vNew(Printer, pwind, rname); 

object  pwind;  /*  parent window       */
char    *rname; /*  report name         */
object  pntr;   /*  new printer object  */
@end group
@end example
This class method is used to create a new @code{Printer} object which
will be used by the @code{Printer} instance methods to control the
printer, fonts and content of output.  This method opens up the
default printer.

@code{pwind} is the window which will act as the parent to any
messages which may be displayed when a report is being printed.
@code{rname} determines what the report is referred to as in any
messages.

The value returned represents the printer to be used.  If the default
printer couldn't be opened, @code{NULL} will be returned.  If an object
is returned, it must be disposed (via @code{Dispose}) when it is no
longer needed.
@example
@group
@exdent Example:

object  pntr, pwind;

pntr = vNew(Printer, pwind, "My Report");
@end group
@end example
@sp 1
See also:  @code{QueryPrinter, NewWithHDC}
@end deffn











@deffn {NewPage} NewPage::Printer
@sp 2
@example
@group
r = gNewPage(pntr); 

object  pntr;   /*  printer object  */
object  r;      /*  printer object  */
@end group
@end example
This method is used to flush and eject the current output page (if anything
had been output to it) and prepare for a possible new page of output.

See the note at the beginning of this section regarding the return value.
@example
@group
@exdent Example:

object  pntr;

if (!gNewPage(pntr))
        abort report;
@end group
@end example
@c @sp 1
@c See also:  @code{}
@end deffn













@deffn {NewWithHDC} NewWithHDC::Printer
@sp 2
@example
@group
pntr = gNewWithHDC(Printer, pwind, rname, hdc); 

object  pwind;  /*  parent window       */
char    *rname; /*  report name         */
HDC     hdc;    /*  device context      */
object  pntr;   /*  new printer object  */
@end group
@end example
This class method is used to create a new @code{Printer} object which
will be used by the @code{Printer} instance methods to control the
printer, fonts and content of output.

This method opens up the printer identified by @code{hdc}.  This is
a Windows internal identifier and may be obtained via the
@code{PrintDialog} class.

@code{pwind} is the window which will act as the parent to any
messages which may be displayed when a report is being printed.
@code{rname} determines what the report is referred to as in any
messages.

The value returned represents the printer to be used.  If the 
printer couldn't be opened, @code{NULL} will be returned.

This method is seldom needed due to @code{QueryPrinter} and @code{New}.
Use them. If an object is returned, it must be disposed (via
@code{Dispose}) when it is no longer needed.
@example
@group
@exdent Example:

object  pntr, pwind;
HDC     hdc;

pntr = gNewWithHDC(Printer, pwind, "My Report", hdc);
@end group
@end example
@sp 1
See also:  @code{QueryPrinter, New}
@end deffn




@deffn {OpenDefault} OpenDefault::Printer
@sp 2
@example
@group
prt = gOpenDefault(Printer, wind, name);

object  wind;  /*  parent window              */
char    *name; /*  report name or NULL        */
object  prt;   /*  new printer object or NULL */
@end group
@end example
This class method is used to open a printer using the system default
printer without showing a print dialog.  @code{wind} is the parent
window for any messages that may be displayed.  @code{name} is the
report name and may be @code{NULL}.

The value returned is a new @code{Printer} instance, or @code{NULL}
if the default printer could not be opened.
@example
@group
@exdent Example:

object  prt, wind;

prt = gOpenDefault(Printer, wind, "My Report");
@end group
@end example
@sp 1
See also:  @code{New, QueryPrinter}
@end deffn



















@deffn {Pie} Pie::Printer
@sp 2
@example
@group
r = gPie(pntr, yBeg, xBeg, yEnd, xEnd, yaBeg, xaBeg, yaEnd, xaEnd);

object  pntr;   /*  printer object         */
int     yBeg;   /*  starting row           */
int     xBeg;   /*  starting column        */
int     yEnd;   /*  ending row             */
int     xEnd;   /*  ending column          */
int     yaBeg;  /*  starting arc row       */
int     xaBeg;  /*  starting arc column    */
int     yaEnd;  /*  ending arc row         */
int     xaEnd;  /*  ending arc column      */
object  r;      /*  printer object         */
@end group
@end example
This method is used to output a pie shaped wedge to the printer.  The
currently selected pen object will indicate the thickness and pattern of
the outline of the shape, and the currently selected brush object will
be used to determine what pattern the shape will be filled with.  The
parameters indicate location of the beginning and ending of the shape as
well as the coordinates of the intersecting line and are in increments
dictated by @code{SetScale}.

See the note at the beginning of this section regarding the return value.
@example
@group
@exdent Example:

object  pntr;

if (!gPie(pntr, 10, 10, 180, 140, 12, 12, 128, 135))
        abort report;
@end group
@end example
@sp 1
See also:  @code{NewPage, SetScale, Use,}
        @code{Line, Rectangle, Ellipse, RoundRect,}
@iftex
@hfil @break @hglue .63in 
@end iftex
@code{Chord, Arc}
@end deffn




@deffn {PrintBitmap} PrintBitmap::Printer
@sp 2
@example
@group
r = gPrintBitmap(prt, file, yPos, xPos, height, width);

object  prt;     /*  printer object          */
char    *file;   /*  bitmap file path        */
int     yPos;    /*  vertical position       */
int     xPos;    /*  horizontal position     */
int     height;  /*  height of output image  */
int     width;   /*  width of output image   */
object  r;       /*  printer object          */
@end group
@end example
This instance method is used to print a bitmap file at the specified
position and size on the current page.  The coordinates are scaled to
printer units.
@example
@group
@exdent Example:

object  prt;

gPrintBitmap(prt, "logo.bmp", 10, 20, 100, 200);
@end group
@end example
@end deffn
















@deffn {Printf} Printf::Printer
@sp 2
@example
@group
r = vPrintf(pntr, fmt, ...);

object   pntr;  /*  printer object     */
char    *fmt;   /*  format string      */
int      r;     /*  length of output   */
@end group
@end example
This method is used to output a string of text (@code{str}) in a sequential
fashion on printer @code{pntr}.  It is analogous to the standard C library
function @code{fprintf}.  All of the standard features of your C library
@code{fprintf} function are supported and that documentation should be
consulted for full documentation on the arguments.

This method is used to support the standard streams interface, is
actually defined by the @code{Stream} class, and is documented here for
convenience.

Note that this method begins with ``v'' since it takes variable arguments.

The value returned by this method is very significant.  During normal
operation, the length of the output string will be returned.  However,
if an error occurs or the user aborted the report a @code{-1} will be
returned.  In this case further output should be avoided (although it
won't hurt - it'll just be ignored) and the printer object should be
disposed.
@example
@group
@exdent Example:

object  pntr;
int     age = 32;

if (-1 == vPrintf(pntr, "My age =- %d\n", age))
        goto abort report;
@end group
@end example
@sp 1
See also:  @code{TextOut, Puts}
@end deffn














@deffn {Puts} Puts::Printer
@sp 2
@example
@group
r = gPuts(pntr, str);

object   pntr;  /*  printer object        */
char    *str;   /*  string to be output   */
int      r;     /*  length of str         */
@end group
@end example
This method is used to output a string of text (@code{str}) in a sequential
fashion on printer @code{pntr}.  This method is used to support the standard
streams interface, is actually defined by the @code{Stream} class, and
is documented here for convenience. @code{str} is the text to be output. 

The value returned by this method is very significant.  During normal
operation, the length of the output string will be returned.  However,
if an error occurs or the user aborted the report a @code{-1} will be
returned.  In this case further output should be avoided (although it
won't hurt - it'll just be ignored) and the printer object should be
disposed.
@example
@group
@exdent Example:

object  pntr;

if (-1 == gPuts(pntr, "Hello World\n"))
        goto abort report;
@end group
@end example
@sp 1
See also:  @code{TextOut, Printf}
@end deffn











@deffn {QueryPrinter} QueryPrinter::Printer
@sp 2
@example
@group
pntr = gQueryPrinter(Printer, pwind, rname); 

object  pwind;  /*  parent window       */
char    *rname; /*  report name         */
object  pntr;   /*  new printer object  */
@end group
@end example
This class method is used to create a new @code{Printer} object which
will be used by the @code{Printer} instance methods to control the
printer, fonts and content of output.

This method queries the user for printer selection and optional
configuration, and uses that information to open the selected
printer.

@code{pwind} is the window which will act as the parent to any
messages which may be displayed when a report is being printed.
@code{rname} determines what the report is referred to as in any
messages.

The value returned represents the printer to be used.  If the user
canceled the printer selection or the printer couldn't be opened,
@code{NULL} will be returned. If an object is returned, it must be
disposed (via @code{Dispose}) when it is no longer needed.
@example
@group
@exdent Example:

object  pntr, pwind;

pntr = gQueryPrinter(Printer, pwind, "My Report");
@end group
@end example
@sp 1
See also:  @code{New, NewWithHDC}
@end deffn
















@deffn {Rectangle} Rectangle::Printer
@sp 2
@example
@group
r = gRectangle(pntr, yBeg, xBeg, yEnd, xEnd);

object  pntr;   /*  printer object   */
int     yBeg;   /*  starting row     */
int     xBeg;   /*  starting column  */
int     yEnd;   /*  ending row       */
int     xEnd;   /*  ending column    */
object  r;      /*  printer object   */
@end group
@end example
This method is used to output a rectangle to the printer.  The currently
selected pen object will indicate the thickness and pattern of the outline
of the shape, and the currently selected brush object will be used to determine
what pattern the shape will be filled with.
The parameters indicate location of the beginning and ending of the shape
and are in increments dictated by @code{SetScale}.

See the note at the beginning of this section regarding the return value.
@example
@group
@exdent Example:

object  pntr;

if (!gRectangle(pntr, 10, 20, 30, 40))
        abort report;
@end group
@end example
@sp 1
See also:  @code{NewPage, SetScale, Use,}
        @code{Line, Ellipse, RoundRect, Chord, Pie, Arc}
@end deffn















@deffn {RoundRect} RoundRect::Printer
@sp 2
@example
@group
r = gRoundRect(pntr, yBeg, xBeg, yEnd, xEnd, eHt, eWth);

object  pntr;   /*  printer object         */
int     yBeg;   /*  starting row           */
int     xBeg;   /*  starting column        */
int     yEnd;   /*  ending row             */
int     xEnd;   /*  ending column          */
int     eHt;    /*  corner ellipse height  */
int     eWth;   /*  corner ellipse width   */
object  r;      /*  printer object         */
@end group
@end example
This method is used to output a rectangle with rounded corners to the
printer.  The currently selected pen object will indicate the thickness
and pattern of the outline of the shape, and the currently selected
brush object will be used to determine what pattern the shape will be
filled with.  The parameters indicate location of the beginning and
ending of the shape and are in increments dictated by @code{SetScale}.

See the note at the beginning of this section regarding the return value.
@example
@group
@exdent Example:

object  pntr;

if (!gRoundRect(pntr, 10, 20, 30, 40, 1, 1))
        abort report;
@end group
@end example
@sp 1
See also:  @code{NewPage, SetScale, Use,}
        @code{Line, Rectangle, Ellipse, Chord, Pie, Arc}
@end deffn




@deffn {Scale} Scale::Printer
@sp 2
@example
@group
r = gScale(prt, row, col);

object  prt;   /*  printer object              */
int     *row;  /*  pointer to row coordinate   */
int     *col;  /*  pointer to column coordinate */
object  r;     /*  printer object              */
@end group
@end example
This instance method scales logical coordinates to physical printer
coordinates in place, accounting for printer size and offsets.
@example
@group
@exdent Example:

object  prt;
int     row = 10, col = 20;

gScale(prt, &row, &col);
@end group
@end example
@end deffn
















@deffn {SetScale} SetScale::Printer
@sp 2
@example
@group
r = gSetScale(pntr, ht, wth);

object  pntr;   /*  printer object  */
int     ht;     /*  logical height  */
int     wth;    /*  logical width   */
object  r;      /*  printer object  */
@end group
@end example
This method is used to set the logical scaling factor associated with an
output device.  Each output page is divided up into @code{ht} evenly
spaced row and @code{wth} evenly spaced columns.  This scaling factor
will be used by all text and graphics output when determining output size
and location on the page.

Coordinates begin at the upper left hand of the page and are zero origin.
The defaults are set to 80 columns and 66 lines.
@example
@group
@exdent Example:

object  pntr;

gSetScale(pntr, 132, 160);
@end group
@end example
@sp 1
See also:  @code{TextOut}
@end deffn












@deffn {TextOut} TextOut::Printer
@sp 2
@example
@group
r = gTextOut(pntr, row, col, txt); 

object  pntr;   /*  printer object  */
int     row;    /*  output row      */
int     col;    /*  output column   */
char    *txt;   /*  text to output  */
object  r;      /*  printer object  */
@end group
@end example
This method is used to output a line of text (@code{txt}) to the printer.
@code{row} and @code{col} indicate the position of the text and are in
increments dictated by @code{SetScale}.  Positions begin in the upper
left hand corner and have an index origin of 0.

The text will be printed in the currently selected font.  Note that the
actual printing will not occur until the report is complete and the
print object disposed.

See the note at the beginning of this section regarding the return value.
@example
@group
@exdent Example:

object  pntr;

if (!gTextOut(pntr, 10, 20, "output text"))
        abort report;
@end group
@end example
@sp 1
See also:  @code{NewPage, Printf, SetScale, LoadFont}
@end deffn







@deffn {Use} Use::Printer
@sp 2
@example
@group
r = gUse(pntr, obj);

object  pntr;   /*  printer object     */
object  obj;    /*  arbitrary object   */
object  r;      /*  the object passed  */
@end group
@end example
This method is used as a general purpose mechanism to associate a
number of object types with a printer.  @code{obj} may be a @code{Font,
Brush} or @code{Pen} object.  Those objects would have normally been
created via their associated classes and then may be associated with a
printer via this method.

Note that the same object should not be associated with more than one
printer object.  The problem is that if one printer object is disposed, WDS
would dispose of all the objects associated with that printer and then the
other printer would reference objects which have been deleted.  The way
around this is to use the @code{Copy} method in order to make a copy of
an object prior to associating it with a new printer.  This way there would
be two independent objects such that if one is deleted the other would
still exist.

The value returned is the object passed.
@example
@group
@exdent Example:

object  myPntr, myOtherPntr, someFont;

someFont = vNew(ExternalFont, "Times New Roman", 12);
if (someFont)  @{
        gUse(myPntr, someFont);
        gUse(myOtherPntr, gCopy(someFont));
@}
@end group
@end example
@sp 1
See also:  @code{LoadFont} and the @code{Brush} and @code{Pen} classes.
@end deffn










@section Menus
The @code{Menu} class, although not used directly by an application, is
used to house the common functionality associated with the
@code{ExternalMenu} and @code{InternalMenu} classes.  Most of the
functionality of those classes is implemented and documented in this
section.

External menus (those defined in resources) would be created and used
via the @code{ExternalMenu} class.  Internal menus (those defined at runtime
by the application code) would be created and used via the
@code{InternalMenu} and @code{PopupMenu} classes.  Since both of these
classes are subclasses of the @code{Menu} class, they inherit most of
their functionality from it.







@deffn {Associate} Associate::Menu
@sp 2
@example
@group
r = mAssociate(menu, id, fun);

object  menu;   /*  a menu object     */
int     id;     /*  menu item id      */
long    (*fun)(object);  /*  function to execute  */
object  r;      /*  the menu object   */
@end group
@end example
This method is used to associate a function with a menu item such that if
the user selects the menu item, @code{fun} gets executed.

In the case of external menus, @code{id} is normally a macro created by
the resource editor while the programmer defines the menu and identifies
a particular menu option.

@code{fun} is the application specific function which automatically gets
executed whenever the user selects the indicated option.  @code{fun} has
the following structure:
@example
@group
static  long    menu_option(object wind)
@{
           .
           .
           .
        return 0L;
@}
@end group
@end example
Where @code{wind} is the window object which the menu is attached to.
The value returned is what is returned by the window procedure associated
with the window.  This value is normally 0 and is fully documented by
the Windows documentation under the @code{WM_COMMAND} message.

The value returned is the menu object passed.
@example
@group
@exdent Example:

object  myMenu;

mAssociate(myMenu, ID_FILE_OPEN, menu_option);
@end group
@end example
@sp 1
See also:  @code{AddMenuOption::InternalMenu, AddMenuOption::PopupMenu}
@end deffn








@deffn {DeepDispose} DeepDispose::Menu
@sp 2
@example
@group
r = gDeepDispose(menu);

object  menu;   /*  a menu object    */
object  r;      /*  NULL             */
@end group
@end example
This method is used to dispose of a menu object and all menu objects
which were pushed (via @code{Push}) into the same menu stack.  It is not
often needed because menus associated with a window are automatically
disposed of (via this method) when the window is disposed.

If @code{menu} is an internal menu, all associated popup menus will
also be disposed.

The value returned is always @code{NULL} and may be used to null out
the variable which contained the object being disposed in order to
avoid future accidental use.
@example
@group
@exdent Example:

object  myMenu;

myMenu = gDeepDispose(myMenu);
@end group
@end example
@sp 1
See also:  @code{Dispose, Push}
@end deffn









@deffn {Dispose} Dispose::Menu
@sp 2
@example
@group
r = gDispose(menu);

object  menu;   /*  a menu object   */
object  r;      /*  NULL            */
@end group
@end example
This method is used to dispose of a menu object when it is no longer
needed.  It is not often needed because menus associated with a window
are automatically disposed of (via @code{DeepDispose}) when the window
is disposed.

If @code{menu} is an internal menu, all associated popup menus will
also be disposed.

The value returned is always @code{NULL} and may be used to null out
the variable which contained the object being disposed in order to
avoid future accidental use.
@example
@group
@exdent Example:

object  myMenu;

myMenu = gDispose(myMenu);
@end group
@end example
@sp 1
See also:  @code{DeepDispose}
@end deffn




@deffn {EnableMenuItem} EnableMenuItem::Menu
@sp 2
@example
@group
prev = gEnableMenuItem(menu, id, flag);

object    menu;   /*  a menu object        */
unsigned  id;     /*  menu item identifier  */
unsigned  flag;   /*  MF_ENABLED, MF_DISABLED, or MF_GRAYED  */
unsigned  prev;   /*  previous state       */
@end group
@end example
This method enables, disables, or grays a menu item identified by
@code{id}.  @code{flag} should be one of @code{MF_ENABLED},
@code{MF_DISABLED}, or @code{MF_GRAYED}.

The value returned is the previous state of the menu item.
@example
@group
@exdent Example:

object  myMenu;

gEnableMenuItem(myMenu, ID_FILE_SAVE, MF_GRAYED);
@end group
@end example
@sp 1
See also:  @code{SetMode}
@end deffn












@deffn {Handle} Handle::Menu
@sp 2
@example
@group
h = gHandle(menu);

object  menu;   /*  menu object    */
HANDLE  h;      /*  Windows handle */
@end group
@end example
This method is used to obtain the Windows internal handle associated with
a menu object.  

Note that this method may be used with most WDS objects in order to obtain
the internal handle that Windows normally associates with each type of object.
See the appropriate documentation.
@example
@group
@exdent Example:

object  myMenu;
HANDLE  h;

h = gHandle(myMenu);
@end group
@end example
@c @sp 1
@c See also:  @code{}
@end deffn









@deffn {MenuFunction} MenuFunction::Menu
@sp 2
@example
@group
fun = mMenuFunction(menu, id);

object  menu;   /*  a menu object     */
int     id;     /*  menu item id      */
long    (*fun)(object);  /*  function to execute  */
@end group
@end example
This method is used to obtain the function which was associated with a
menu item.
@c @example
@c @group
@c @exdent Example:
@c @end group
@c @end example
@sp 1
See also:  @code{Associate}
@end deffn







@deffn {Pop} Pop::Menu
@sp 2
@example
@group
prev = gPop(top);

object  top;    /*  top menu       */
object  prev;   /*  previous menu  */
@end group
@end example
This method is used to dispose of menu @code{top} and gain access to the
next menu in menu @code{top}'s last in, first out list.  It is mainly
used by the @code{Window} class in order to support the ability to
change a menu and later return to a previous menu.  The value returned
is the next menu in the stack.
@example
@group
@exdent Example:

object  myMenu, subMenu;
HANDLE  h;

subMenu = gPop(myMenu);
@end group
@end example
@sp 1
See also:  @code{Push, Use::Window, PopMenu::Window}
@end deffn






@deffn {Push} Push::Menu
@sp 2
@example
@group
r = gPush(top, psh);

object  top;    /*  top menu     */
object  psh;    /*  pushed menu  */
object  r;      /*  top menu     */
@end group
@end example
This method is used to push menu (@code{psh}) into a last in, first out
stack accessible through menu @code{top}.  It is mainly used by the
@code{Window} class in order to support the ability to change a menu and
later return to a previous menu.  The value returned is the top menu
passed.
@example
@group
@exdent Example:

object  myMenu, subMenu;
HANDLE  h;

h = gPush(myMenu, subMenu);
@end group
@end example
@sp 1
See also:  @code{Pop, Use::Window, PopMenu::Window}
@end deffn




@deffn {SetCheckMark} SetCheckMark::Menu
@sp 2
@example
@group
prev = gSetCheckMark(menu, id, val);

object    menu;   /*  a menu object        */
unsigned  id;     /*  menu item identifier  */
int       val;    /*  non-zero to check, zero to uncheck  */
int       prev;   /*  previous state       */
@end group
@end example
This method checks or unchecks a menu item identified by @code{id}.
If @code{val} is non-zero, the menu item is checked; if zero, it is
unchecked.

The value returned is the previous check state.
@example
@group
@exdent Example:

object  myMenu;

gSetCheckMark(myMenu, ID_VIEW_TOOLBAR, 1);
@end group
@end example
@end deffn








@deffn {SetMode} SetMode::Menu
@sp 2
@example
@group
r = mSetMode(menu, id, mode);

object  menu;   /*  menu       */
unsigned  id;   /*  menu item  */
unsigned  mode; /*  item mode  */
object    r;    /*  menu       */
@end group
@end example
This method is used to set the mode associated with a menu item.  This
mode specifies whether the menu is enabled, disabled or grayed.

In the case of external menus, @code{id} is normally a macro created by
the resource editor while the programmer defines the menu and identifies
a particular menu option.  In the case of internal menus, @code{id}
is the number returned by @code{AddMenuOption::InternalMenu} or
@code{AddMenuOption::PopupMenu}.

@code{mode} is a Windows defined macro and may be one of the following:
@example
@group
MF_ENABLED      to enable the menu item
MF_DISABLED     to disable the menu item
MF_GRAYED       to gray the menu option
@end group
@end example
The value returned is the menu passed.
@example
@group
@exdent Example:

object  myMenu;

mSetMode(myMenu, ID_FILE_OPEN, MF_DISABLED);
@end group
@end example
@sp 1
See also:  @code{AddMenuOption::InternalMenu, AddMenuOption::PopupMenu}
@end deffn






@subsection External Menus
The @code{ExternalMenu} class is used to access menus which were defined
using the resource editor.  Since this class is a subclass of @code{Menu},
all of @code{Menu}'s functionality is accessible through instance of this
class.  Therefore, see the @code{Menu} class for additional functionality.







@deffn {Load} Load::ExternalMenu
@sp 2
@example
@group
menu = mLoad(ExternalMenu, id);

unsigned  id;   /*  menu id  */
object  menu;   /*  menu     */
@end group
@end example
This method is used to create a new menu object representing a menu defined
with the resource editor.

@code{id} is normally a macro created by the resource editor which identifies
a particular menu.

The value returned is a new object representing the menu or @code{NULL} if
the menu identified is not found.
@example
@group
@exdent Example:

object  myMenu;

myMenu = mLoad(ExternalMenu, MY_MENU);
@end group
@end example
@sp 1
See also:  @code{LoadStr, LoadMenu::Window, LoadMenuStr::Window}
@end deffn







@deffn {LoadStr} LoadStr::ExternalMenu
@sp 2
@example
@group
menu = gLoadStr(ExternalMenu, id);

char    *id;    /*  menu id  */
object  menu;   /*  menu     */
@end group
@end example
This method is used to create a new menu object representing a menu defined
with the resource editor.

@code{id} is a string which names the desired menu.  This name is
defined by the resource editor and identifies a particular menu.

The value returned is a new object representing the menu or @code{NULL} if
the menu identified is not found.
@example
@group
@exdent Example:

object  myMenu;

myMenu = gLoadStr(ExternalMenu, "mymenu");
@end group
@end example
@sp 1
See also:  @code{Load, LoadMenu::Window, LoadMenuStr::Window}
@end deffn








@subsection Internal Menus
The @code{InternalMenu} class is used to enable an application to create
menus on the fly -- without the need to pre-specify the menu's structure
with the resource editor.  Since this class is a subclass of
@code{Menu}, all of @code{Menu}'s functionality is accessible through
instance of this class.  Therefore, see the @code{Menu} class for
additional functionality.








@deffn {AddMenuOption} AddMenuOption::InternalMenu
@sp 2
@example
@group
id = gAddMenuOption(mnu, ttl, fun)

object  menu;   /*  menu     */
char    *ttl;   /*  title    */
long    (*fun)(); /*  function  */
int     id;     /*  item id  */
@end group
@end example
This method is used to append a new option to internal menu object
@code{mnu} and associate it to an application specific function
(@code{fun}).  Once this is done, if the user selects the new menu
option the function @code{fun} will be executed.

@code{ttl} represents the character string which will be displayed
in the menu.  @code{ttl} may also contain an embedded ampersand (&)
character.  This causes the character following the ampersand to
be underlined.  It also causes the following character to be the
Alt key selection the user can use to quickly select a menu option.

@code{fun} is the function which will be executed when the user selects
the menu option and has the following form:
@example
@group
long    fun(object wind)
@{
        .
        .
        .
        return 0L;
@}
@end group
@end example
The function executed (@code{fun}) is passed the window object and
returns a long.  The return value is documented in the Windows documentation
under the message named @code{WM_COMMAND}.  It should normally be @code{0L}.

This method returns a unique WDS generated id which identifies this
particular menu option and may be used to control the status of the option.
@example
@group
@exdent Example:

static  long    file_message(object wind)
@{
        gMessage(wind, "File_Message");
        return 0L;
@}

        .
        .
        mnu = vNew(InternalMenu);
        gAddMenuOption(mnu, "&File", file_message);
@end group
@end example
@sp 1
See also:  @code{AddPopupMenu}
@end deffn











@deffn {AddPopupMenu} AddPopupMenu::InternalMenu
@sp 2
@example
@group
id = gAddPopupMenu(mnu, ttl, pm)

object  mnu;    /*  menu        */
char    *ttl;   /*  title       */
object  pm;     /*  Popup menu  */
object  r;      /*  menu        */
@end group
@end example
This method is used to append a new option to menu object @code{mnu} and
associate it to a popup menu.  Once this is done, if the user selects
the new menu option the popup menu will be displayed.

@code{ttl} represents the character string which will be displayed
in the menu.  @code{ttl} may also contain an embedded ampersand (&)
character.  This causes the character following the ampersand to
be underlined.  It also causes the following character to be the
Alt key selection the user can use to quickly select a menu option.

@code{pm} is a popup menu object created with the @code{PopupMenu}
class.

This method returns @code{mnu}.
@example
@group
@exdent Example:

object  mnu, pm;

mnu = vNew(InternalMenu);
pm  = vNew(PopupMenu, mnu);
gAddPopupMenu(mnu, "&File", pm);
@end group
@end example
@sp 1
See also:  @code{AddMenuOption, New::PopupMenu}
@end deffn








@deffn {New} New::InternalMenu
@sp 2
@example
@group
menu = vNew(InternalMenu);

object  menu;   /*  menu     */
@end group
@end example
This method is used to create a new internal menu.  An internal menu is
one which is not pre-specified with the resource editor.  It may therefore
be structured at runtime.

The object returned is a new internal menu object with no items associated
with it.
@example
@group
@exdent Example:

object  myMenu;

myMenu = vNew(InternalMenu);
@end group
@end example
@sp 1
See also:  @code{AddMenuOption, AddMenu, Use::Window, LoadMenu::Window}
@end deffn





@section Popup Menus
The @code{PopupMenu} class is used in conjunction with the
@code{InternalMenu} class in order to provide a nested menu structure.
Nesting of popup menus is supported to arbitrary levels.








@deffn {AddMenuOption} AddMenuOption::PopupMenu
@sp 2
@example
@group
id = gAddMenuOption(mnu, ttl, fun)

object  menu;   /*  menu     */
char    *ttl;   /*  title    */
long    (*fun)(); /*  function  */
int     id;     /*  item id  */
@end group
@end example
This method works exactly like @code{AddMenuOption::InternalMenu}.
See that description for full documentation.
@c @example
@c @group
@c @exdent Example:
@c @end group
@c @end example
@sp 1
See also:  @code{AddPopupMenu, AddSeparator}
@end deffn









@deffn {AddPopupMenu} AddPopupMenu::PopupMenu
@sp 2
@example
@group
id = gAddPopupMenu(mnu, ttl, pm)

object  mnu;    /*  menu        */
char    *ttl;   /*  title       */
object  pm;     /*  Popup menu  */
object  r;      /*  menu        */
@end group
@end example
@code{mnu} is the popup menu which is to act as the parent of @code{pm}
and may be arbitrarily nested.

This method is the same as @code{AddPopupMenu::InternalMenu}.  See that
method for complete documentation.
@c @example
@c @group
@c @exdent Example:
@c @end group
@c @end example
@sp 1
See also:  @code{AddMenuOption, AddSeparator}
@end deffn










@deffn {AddSeparator} AddSeparator::PopupMenu
@sp 2
@example
@group
r = gAddSeparator(mnu)

object  mnu;    /*  menu        */
object  r;      /*  menu        */
@end group
@end example
This method is used to append a horizontal bar to the end of a popup
menu.  Any menu items added subsequently will appear after the
horizontal separator.

The menu object passed is returned.
@example
@group
@exdent Example:

gAddSeparator(pm);
@end group
@end example
@sp 1
See also:  @code{AddMenuOption}
@end deffn











@deffn {Dispose} Dispose::PopupMenu
@sp 2
@example
@group
r = gDispose(menu);

object  menu;   /*  a menu object   */
object  r;      /*  NULL            */
@end group
@end example
This method is used to dispose of a menu object when it is no longer
needed.  It is not often needed because popup menus associated with
other menus are automatically disposed of (via @code{DeepDispose}) when
the top level menu is disposed.

This class also provides a @code{DeepDispose} which performs the same
function.

The value returned is always @code{NULL} and may be used to null out
the variable which contained the object being disposed in order to
avoid future accidental use.
@example
@group
@exdent Example:

object  myMenu;

myMenu = gDispose(myMenu);
@end group
@end example
@sp 1
See also:  @code{DeepDispose::Menu}
@end deffn








@c Handle::PopupMenu removed — gHandle is not implemented in PopupMenu.








@deffn {MenuFunction} MenuFunction::PopupMenu
@sp 2
@example
@group
fun = mMenuFunction(menu, id);

object  menu;   /*  a menu object     */
int     id;     /*  menu item id      */
long    (*fun)(object);  /*  function to execute  */
@end group
@end example
This method is used to obtain the function which was associated with a
menu item.
@c @example
@c @group
@c @exdent Example:
@c @end group
@c @end example
@sp 1
See also:  @code{AddMenuOption}
@end deffn








@deffn {New} New::PopupMenu
@sp 2
@example
@group
pm = vNew(PopupMenu, tm);

object  tm;     /*  top menu    */
object  pm;     /*  popup menu  */
@end group
@end example
This method is used to create a new popup menu.  A popup menu is one
which pops up or appears when a user selects an associated menu option.
Popup menus are used in conjunction with internal menus in order to
create arbitrarily complex menu structures.

@code{tm} is the top menu object.  It must be an @code{InternalMenu},
regardless of where the new popup menu will be nested.

The object returned represents the new popup menu created and may be
used to add items to that menu.
@example
@group
@exdent Example:

object  myMenu, pm;

myMenu = vNew(InternalMenu);
pm = vNew(PopupMenu, myMenu);
gAddPopupMenu(myMenu, "&File", pm);
@end group
@end example
@sp 1
See also:  @code{AddMenuOption, AddPopupMenu, AddSeparator,}
@iftex
@hfil @break @hglue .63in     
@end iftex
@code{AddPopupMenu::InternalMenu}
@end deffn







